% ch7_b (书页 287-298 / p302-p313)
%--A-TAIL--
%JOIN%尘埃的静止系，其中能量-动量张量取如下形式
\begin{equation*}
T_{\mu\nu} = \rho U_\mu U_\nu =
\begin{pmatrix}
\rho & & & \\
& 0 & & \\
& & 0 & \\
& & & 0
\end{pmatrix}.
\tag{7.54}
\end{equation*}

% ---- 书页 287 / p302 ----（页顶承接句与 (7.54) 已在上面 tail 中）

由于我们的背景是平直 Minkowski 时空，要容纳运动的源并不困难，只需通过洛伦兹变换换到其静止系即可；在这个极限下我们无法处理的，是具有较大相对速度的多个源。

现在转向横向规范下的 Einstein 方程，即式 (7.38)--(7.40)。对于静态源，我们去掉所有含时间导数的项，同时代入能量-动量张量 (7.54)，得到
\begin{align*}
\nabla^2\Psi &= 4\pi G\rho \\
\nabla^2 w_j &= 0 \\
(\delta_{ij}\nabla^2 - \partial_i\partial_j)(\Phi - \Psi) - \nabla^2 s_{ij} &= 0.
\tag{7.55}
\end{align*}
我们将寻找既无奇点、在无穷远处又表现良好的解；因此，只有那些由右端项作为源的场才会非零。例如，(7.55) 中的第二个方程直接给出 $w^i = 0$。接下来对第三个方程取迹（对 $\delta^{ij}$ 求和）：
\begin{equation*}
2\nabla^2(\Phi - \Psi) = 0.
\tag{7.56}
\end{equation*}
这就迫使两个标量势相等，
\begin{equation*}
\Phi = \Psi.
\tag{7.57}
\end{equation*}
回想在第 4 章最初讨论 Newton 极限时，我们论证过微扰的 00 分量 $\Phi$（正是它决定了非相对论性粒子的运动）满足 Poisson 方程；而从 (7.55) 来看，满足这个方程的似乎其实是空间分量上的标量微扰 $\Psi$。二者之间隐含的联系由 (7.56) 给出：当 $T_{ij}$ 的迹（三个主压强之和）为零时，它令两个势相等。最后，把 $\Phi = \Psi$ 代入 (7.55) 的最后一个方程，得到
\begin{equation*}
\nabla^2 s_{ij} = 0,
\tag{7.58}
\end{equation*}
对于表现良好的解，这意味着 $s_{ij} = 0$。

因此，静态 Newton 源的微扰度规为
\begin{equation*}
\boxed{\,ds^2 = -(1+2\Phi)\mathrm{d}t^2 + (1-2\Phi)(\mathrm{d}x^2 + \mathrm{d}y^2 + \mathrm{d}z^2),\,}
\tag{7.59}
\end{equation*}

% ---- 书页 288 / p303 ----

%FIG{7.3}: {一条被偏折的测地线 $x^{\mu}(\lambda)$，分解为一条背景测地线 $x^{(0)\mu}$ 和一个微扰 $x^{(1)\mu}$。偏转角 $\widehat{\alpha}$ 代表波矢量（wave vector）沿路径转过的量的（负值）。图中画出的是碰撞参数（impact parameter）为 $b$ 的单个质量 $M$，尽管这一设置更具一般性。}

或等价地
\begin{equation*}
h_{\mu\nu} =
\begin{pmatrix}
-2\Phi & & & \\
& -2\Phi & & \\
& & -2\Phi & \\
& & & -2\Phi
\end{pmatrix},
\tag{7.60}
\end{equation*}
其中势满足通常的 Poisson 方程，
\begin{equation*}
\nabla^2\Phi = 4\pi G\rho.
\tag{7.61}
\end{equation*}
这是对第 4 章结果的一个重要推广，因为我们现在不仅知道 $h_{00}$，还知道空间度规的微扰。

现在让我们考虑光子（或其他无质量粒子）在这个几何中的路径；换句话说，对类光轨迹 $x^\mu(\lambda)$ 求解微扰后的测地线方程。\footnote{我们所用的方法在 T.~Pyne 和 M.~Birkinshaw 的文章 \emph{Astrophys. Journ.} 458, 46 (1996) 中有概述，见 \texttt{http://arxiv.org/abs/astro-ph/9504060}。}我们所考虑的几何描绘于图 7.3。回想一下，我们的处理思路是把度规微扰当作定义在平直背景时空上的场。类似地，我们可以把测地线分解为背景路径加上微扰，
\begin{equation*}
x^\mu(\lambda) = x^{(0)\mu}(\lambda) + x^{(1)\mu}(\lambda),
\tag{7.62}
\end{equation*}
其中 $x^{(0)\mu}$ 满足背景中的测地线方程（换句话说，就是一条笔直的类光路径）。然后我们沿背景路径取所有量的值，来求解 $x^{(1)\mu}(\lambda)$。要让这一做法行得通，我们需要假设势 $\Phi$ 沿背景测地线与真实测地线没有明显差别；这个条件相当于要求 $x^{(1)i}\partial_i\Phi \ll \Phi$。不过，即使这个条件不成立，也并非全盘皆输。如果我们只考虑非常短的路径，偏离

% ---- 书页 289 / p304 ----

$x^{(1)\mu}$ 就必然很小，我们的近似也将有效。而接下来我们可以用这样的短段拼接出更长的路径。结果是，我们会推导出真实的方程，只是积分所沿的路径将是\emph{真实}路径 $x^\mu(\lambda)$，而不是背景路径 $x^{(0)\mu}(\lambda)$。只要理解了这一点，我们的结果对微扰时空中的任何轨迹都将成立。

为方便起见，我们把背景路径的波矢量记作 $k^\mu$，把偏离矢量的导数记作 $\ell^\mu$：
\begin{equation*}
k^\mu \equiv \frac{dx^{(0)\mu}}{d\lambda}, \qquad \ell^\mu \equiv \frac{dx^{(1)\mu}}{d\lambda}.
\tag{7.63}
\end{equation*}
路径为类光的条件当然是
\begin{equation*}
g_{\mu\nu}\frac{dx^\mu}{d\lambda}\frac{dx^\nu}{d\lambda} = 0,
\tag{7.64}
\end{equation*}
我们必须逐阶求解它。零阶时我们直接有 $\eta_{\mu\nu}k^\mu k^\nu = 0$，即
\begin{equation*}
(k^0)^2 = (\vec{k}\,)^2 \equiv k^2,
\tag{7.65}
\end{equation*}
其中 $\vec{k}$ 是分量为 $k^i$ 的三维矢量。这个方程充当常数 $k$ 的定义。于一阶我们得到
\begin{equation*}
2\eta_{\mu\nu}k^\mu\ell^\nu + h_{\mu\nu}k^\mu k^\nu = 0,
\tag{7.66}
\end{equation*}
即
\begin{equation*}
-k\ell^0 + \vec{\ell}\cdot\vec{k} = 2k^2\Phi.
\tag{7.67}
\end{equation*}

现在我们转向微扰后的测地线方程，
\begin{equation*}
\frac{d^2x^\mu}{d\lambda^2} + \Gamma^\mu_{\rho\sigma}\frac{dx^\rho}{d\lambda}\frac{dx^\sigma}{d\lambda} = 0.
\tag{7.68}
\end{equation*}
把 $w^i = 0$ 和 $h_{ij} = -2\Phi\delta_{ij}$ 代入 (7.19)，便可以求出 Christoffel 符号：
\begin{align*}
\Gamma^0_{0i} &= \Gamma^i_{00} = \partial_i\Phi, \\
\Gamma^i_{jk} &= \delta_{jk}\partial_i\Phi - \delta_{ik}\partial_j\Phi - \delta_{ij}\partial_k\Phi.
\tag{7.69}
\end{align*}
零阶测地线方程只是告诉我们 $x^{(0)\mu}$ 是一条直线轨迹，而在一阶我们有
\begin{equation*}
\frac{d\ell^\mu}{d\lambda} = -\Gamma^\mu_{\rho\sigma}k^\rho k^\sigma.
\tag{7.70}
\end{equation*}

% ---- 书页 290 / p305 ----

右端没有 $\ell^\mu$ 的因子，因为 Christoffel 符号已经是微扰的一阶量。(7.70) 的 $\mu = 0$ 分量为
\begin{equation*}
\frac{d\ell^0}{d\lambda} = -2k(\vec{k}\cdot\vec{\nabla}\Phi),
\tag{7.71}
\end{equation*}
而空间分量为
\begin{equation*}
\frac{d\vec{\ell}}{d\lambda} = -2k^2\vec{\nabla}_{\perp}\Phi.
\tag{7.72}
\end{equation*}
这里我们引入了横向于路径的梯度，它定义为总梯度减去沿路径方向的梯度，
\begin{align*}
\vec{\nabla}_{\perp}\Phi &\equiv \vec{\nabla}\Phi - \vec{\nabla}_{\parallel}\Phi \\
&= \vec{\nabla}\Phi - k^{-2}(\vec{k}\cdot\vec{\nabla}\Phi)\vec{k}.
\tag{7.73}
\end{align*}
在以上所有表达式中，路径均指背景路径。

注意，在 $\Phi$ 的一阶精度下，空间波矢量微扰 $\vec{\ell}$ 与原有的空间波矢量 $\vec{k}$ 正交。为看清这一点，我们可以积分 (7.71) 得到 $\ell^0$ 的表达式：
\begin{align*}
\ell^0 &= \int \frac{d\ell^0}{d\lambda}\, d\lambda \\
&= -2k\int (\vec{k}\cdot\vec{\nabla}\Phi)\, d\lambda \\
&= -2k\int \left(\frac{d\vec{x}}{d\lambda}\cdot\vec{\nabla}\Phi\right) d\lambda \\
&= -2k\int \vec{\nabla}\Phi\cdot d\vec{x} \\
&= -2k\Phi.
\tag{7.74}
\end{align*}
积分常数由 $\Phi = 0$ 时 $\ell^0 = 0$ 的要求确定。把它代入 (7.67)，得到
\begin{equation*}
\vec{\ell}\cdot\vec{k} = k\ell^0 + 2k^2\Phi = 0,
\tag{7.75}
\end{equation*}
这就验证了 $\vec{\ell}$ 与 $\vec{k}$ 在一阶精度下正交。

\textbf{偏转角}（deflection angle）$\widehat{\alpha}$ 是原有的空间波矢量从源传播到观者处时被偏折的量；它是垂直于 $\vec{k}$ 的平面内的一个二维矢量。（我们采用记号 $\widehat{\alpha}$ 而不是 $\vec{\alpha}$，因为后者已留给第 8 章引入的约化偏转角（reduced deflection angle）。）由图 7.3 所描绘的几何，偏转角可以表示为

% ---- 书页 291 / p306 ----

\begin{equation*}
\widehat{\alpha} = -\frac{\Delta\vec{\ell}}{k},
\tag{7.76}
\end{equation*}
其中负号只是考虑到这样一个事实：偏转角是由沿光子路径向后看的观者测量的。波矢量的转动可以由 (7.72) 算出：
\begin{align*}
\Delta\vec{\ell} &= \int\frac{d\vec{\ell}}{d\lambda}\, d\lambda \\
&= -2k^2\int \vec{\nabla}_{\perp}\Phi\, d\lambda.
\tag{7.77}
\end{align*}
因此，偏转角可以写成沿所经过的物理空间距离 $s = k\lambda$ 的积分：
\begin{equation*}
\boxed{\,\widehat{\alpha} = 2\int \vec{\nabla}_{\perp}\Phi\, ds.\,}
\tag{7.78}
\end{equation*}

我们可以在点质量的情形下算出偏转角：设想背景路径沿 $x$ 方向，碰撞参数由一个横向矢量 $\vec{b}$ 定义，它在最接近点处从路径指向质量。取 $b = |\vec{b}|$，势为
\begin{equation*}
\Phi = -\frac{GM}{r} = -\frac{GM}{(b^2+x^2)^{1/2}},
\tag{7.79}
\end{equation*}
其横向梯度因此为
\begin{equation*}
\vec{\nabla}_{\perp}\Phi = \frac{GM}{(b^2+x^2)^{3/2}}\vec{b}.
\tag{7.80}
\end{equation*}
于是偏转角为
\begin{align*}
\widehat{\alpha} &= 2GMb\int \frac{dx}{(b^2+x^2)^{3/2}} \\
&= \frac{4GM}{b},
\tag{7.81}
\end{align*}
其中积分取自 $-\infty$ 到 $\infty$，这里假定了源和观者都离偏折质量非常远。注意在我们的单位制中 $c = 1$；在其他单位制下，应在 (7.81) 的分母中补上因子 $c^2$。

太阳对光的偏折在历史上是广义相对论的一个关键检验。Einstein 提出了三项这样的检验：水星近日点进动、引力红移和光线偏折。水星近日点进动被 GR 成功解释，但它解释的是一个早已被观测到的偏差；引力红移直到很久之后才被观测到，
% ---- 书页 292 / p307 ----
因此光线偏折是 Einstein 的理论第一次正确预言了尚未被探测到的现象。Eddington 率领的一支著名远征队在 1919 年的日全食期间观测了太阳附近恒星的位置；观测结果与 GR 预言相符，世界各地的报纸都在头版报道了这一消息。预言的效应相当小：对太阳而言，$GM_\odot/c^2 = 1.48\times10^5$ cm，太阳半径为 $R_\odot = 6.96\times10^{10}$ cm，由此给出的最大偏转角为 $\widehat{\alpha} = 1.75$ 角秒。后来对 Eddington 结果的重新评估使人们怀疑他是否真的达到了当初宣称的精度；当代的测量则利用从太阳背后经过的类星体（quasar）的高精度干涉观测，对 GR 作出非常精确的检验（GR 迄今为止都经受了这些检验）。与此同时，对星系和恒星等天体物理源造成的光线偏折的观测，已经发展成一个活跃的研究领域，其名目叫作“引力透镜”。当然，在这种情形下我们对透镜质量的了解很少能达到为 GR 提供精密检验的程度；更常见的做法反而是把观测到的偏转角用作测量质量的一种手段。我们将在第 8 章进一步讨论透镜。

除了光线偏折之外，1964 年 Shapiro 还指出了弱场广义相对论在光子轨迹上的另一个可观测结果：引力时间延迟（gravitational time delay）。沿一条类光路径所耗费的坐标时间为
\begin{equation*}
t = \int\frac{dx^0}{d\lambda}\, d\lambda.
\tag{7.82}
\end{equation*}
我们把自己放在一个远离一切源、静止于背景惯性系中的观者的位置上，因此坐标时间就是我们的固有时。在存在 Newton 势的情况下，光子相对于背景光锥看起来“慢了下来”，由此导致的附加时间延迟为
\begin{align*}
\Delta t &= \int\frac{dx^{(1)0}}{d\lambda}\, d\lambda \\
&= \int \ell^0\, d\lambda \\
&= -2k\int \Phi\, d\lambda,
\tag{7.83}
\end{align*}
即
\begin{equation*}
\Delta t = -2\int \Phi\, ds.
\tag{7.84}
\end{equation*}

按照我们的规则，积分沿背景路径进行。除了这个 Shapiro 延迟之外，还可能有一个额外的“几何”时间延迟，因为真实路径经过的空间距离比背景路径更长。对太阳造成的光线偏折而言，几何延迟效应可以忽略，但在宇宙学的应用中，
% ---- 书页 293 / p308 ----
它可以与 Shapiro 效应相比拟。这一时间延迟已被观测到，其中最精确的测量利用的是航天器而非天然天体；详见 Will (1981)。

光子穿过 Newton 势的运动——它同时导致光线偏折和引力时间延迟——也可以等价地这样导出：设想光子是在折射率（refractive index）为
\begin{equation*}
n = 1 - 2\Phi,
\tag{7.85}
\end{equation*}
的介质中传播（精确到一阶）。事实上，我们本可以利用 Fermat 最短时间原理（Fermat's principle of least time）求出光子的运动方程；习题中将请你证明这一点。

\section{引力波解}

弱场极限一个更令人兴奋的应用是引力辐射。这里我们研究的是引力场的自由传播自由度，它们的存在不需要任何局域的源（当然它们可以由这类源产生）。于是我们再一次回到横向规范下的弱场方程，即式 (7.38)--(7.40)，但这一次保留时间导数，同时完全关掉能量-动量张量，$T_{\mu\nu} = 0$。这时 00 方程为
\begin{equation*}
\nabla^2\Psi = 0,
\tag{7.86}
\end{equation*}
在表现良好的边界条件下，这意味着 $\Psi = 0$。于是 $0j$ 方程为
\begin{equation*}
\nabla^2 w_j = 0,
\tag{7.87}
\end{equation*}
这同样意味着 $w_j = 0$。

接下来看 $ij$ 方程的迹，（代入上面的结果后）它给出
\begin{equation*}
\nabla^2\Phi = 0,
\tag{7.88}
\end{equation*}
这意味着 $\Phi = 0$。

于是只剩下 $ij$ 方程的无迹部分，它成为无迹应变张量（strain tensor）的波动方程：
\begin{equation*}
\Box s_{ij} = 0.
\tag{7.89}
\end{equation*}
虽然到目前为止用 $s_{ij}$ 做事很方便，但文献中远为常见的做法是用整个度规微扰 $h_{\mu\nu}$ 来书写表达式，不过采用的拟设是令所有其他自由度 $(\Phi, \Psi, w_i)$ 都为零（且 $s_{ij}$ 是横向的）。这就是通常所说的\textbf{横无迹（TT）规范}（transverse-traceless gauge），其中有

% ---- 书页 294 / p309 ----

\begin{equation*}
h^{\mathrm{TT}}_{\mu\nu} =
\begin{pmatrix}
0 & 0 & 0 & 0 \\
0 & & & \\
0 & & 2s_{ij} & \\
0 & & &
\end{pmatrix}.
\tag{7.90}
\end{equation*}
这时的运动方程为
\begin{equation*}
\Box h^{\mathrm{TT}}_{\mu\nu} = 0.
\tag{7.91}
\end{equation*}
为了更容易与其他资料对照，在讨论引力波时我们将使用 $h^{\mathrm{TT}}_{\mu\nu}$ 而不是 $s_{ij}$，同时记住 $h^{\mathrm{TT}}_{\mu\nu}$ 是纯空间的、无迹且横向的：
\begin{gather*}
h^{\mathrm{TT}}_{0\nu} = 0 \\
\eta^{\mu\nu}h^{\mathrm{TT}}_{\mu\nu} = 0 \\
\partial_\mu h^{\mu\nu}_{\mathrm{TT}} = 0.
\tag{7.92}
\end{gather*}

我们从波动方程 (7.91) 出发寻找解。熟悉电磁学中类似问题的读者会注意到，这里的步骤几乎一模一样。这族波动方程中特别有用的一组解是平面波，由下式给出
\begin{equation*}
h^{\mathrm{TT}}_{\mu\nu} = C_{\mu\nu}e^{ik_\sigma x^\sigma},
\tag{7.93}
\end{equation*}
其中 $C_{\mu\nu}$ 是一个常值、对称的 $(0,2)$ 型张量，显然无迹且纯空间：
\begin{gather*}
C_{0\nu} = 0 \\
\eta^{\mu\nu}C_{\mu\nu} = 0.
\tag{7.94}
\end{gather*}
当然，$e^{ik_\sigma x^\sigma}$ 是复的，而 $h^{\mathrm{TT}}_{\mu\nu}$ 是实的；我们让实部和虚部一起贯穿整个计算，最后再取实部。常矢量 $k^\sigma$ 就是波矢量。为了确认我们得到的确是一个解，把它代入：
\begin{align*}
0 &= \Box h^{\mathrm{TT}}_{\mu\nu} \\
&= \eta^{\rho\sigma}\partial_\rho\partial_\sigma h^{\mathrm{TT}}_{\mu\nu} \\
&= \eta^{\rho\sigma}\partial_\rho(ik_\sigma h^{\mathrm{TT}}_{\mu\nu}) \\
&= -\eta^{\rho\sigma}k_\rho k_\sigma h^{\mathrm{TT}}_{\mu\nu} \\
&= -k_\sigma k^\sigma h^{\mathrm{TT}}_{\mu\nu}.
\tag{7.95}
\end{align*}

% ---- 书页 295 / p310 ----

由于对一个有意义的解来说，$h^{\mathrm{TT}}_{\mu\nu}$ 的各分量不会处处都为零，我们必须有
\begin{equation*}
k_\sigma k^\sigma = 0.
\tag{7.96}
\end{equation*}
因此，只要波矢量是类光的，平面波 (7.93) 就是线性化方程的解；这可以粗略地转述为：引力波以光速传播。波矢量的类时分量就是波的频率，我们把它写作 $k^\sigma = (\omega, k^1, k^2, k^3)$。（更一般地，以四速度 $U^\mu$ 运动的观者会观测到波的频率为 $\omega = -k_\mu U^\mu$。）于是波矢量为类光的条件就成为
\begin{equation*}
\omega^2 = \delta_{ij}k^ik^j.
\tag{7.97}
\end{equation*}

当然，我们的波远非最一般的解；任意（可能无限多个）不同的平面波可以相加，而且仍然求解线性方程 (7.91)。事实上，任何解都可以写成这样的叠加。

我们还需要保证微扰是横向的。这意味着
\begin{align*}
0 &= \partial_\mu h^{\mu\nu}_{\mathrm{TT}} \\
&= iC^{\mu\nu}k_\mu e^{ik_\sigma x^\sigma},
\tag{7.98}
\end{align*}
而这只有当
\begin{equation*}
k_\mu C^{\mu\nu} = 0.
\tag{7.99}
\end{equation*}
时才成立。我们说波矢量与 $C^{\mu\nu}$ 正交。

选取空间坐标使波沿 $x^3$ 方向传播，可以把我们的解写得更显式；也就是说，
\begin{equation*}
k^\mu = (\omega, 0, 0, k^3) = (\omega, 0, 0, \omega),
\tag{7.100}
\end{equation*}
其中我们知道 $k^3 = \omega$，因为波矢量是类光的。在这种情形下，$k^\mu C_{\mu\nu} = 0$ 与 $C_{0\nu} = 0$ 合起来意味着
\begin{equation*}
C_{3\nu} = 0.
\tag{7.101}
\end{equation*}
于是 $C_{\mu\nu}$ 仅有的非零分量是 $C_{11}$、$C_{12}$、$C_{21}$ 和 $C_{22}$。但 $C_{\mu\nu}$ 无迹且对称，故一般地可以写成
\begin{equation*}
C_{\mu\nu} =
\begin{pmatrix}
0 & 0 & 0 & 0 \\
0 & C_{11} & C_{12} & 0 \\
0 & C_{12} & -C_{11} & 0 \\
0 & 0 & 0 & 0
\end{pmatrix}.
\tag{7.102}
\end{equation*}
这样，对这个规范中沿 $x^3$ 方向传播的平面波，两个分量 $C_{11}$ 和 $C_{12}$（连同频率 $\omega$）就完全刻画了这个波。

% ---- 书页 296 / p311 ----

为了体会引力波经过时的物理效应，考虑检验粒子在波存在时的运动。只求解单个粒子的轨迹肯定是不够的，因为那只会告诉我们世界线上坐标的取值。事实上，对任何一个单独的粒子，我们都能找到一组横无迹坐标（transverse traceless coordinates），使这个粒子在精确到 $h^{\mathrm{TT}}_{\mu\nu}$ 一阶的意义下显得静止。为了得到一个不依赖坐标的量来量度波的效应，我们考虑邻近粒子的相对运动，这由测地偏离方程描述。如果考虑一些邻近的粒子，其四速度由同一个矢量场 $U^\mu(x)$ 描写，间隔矢量为 $S^\mu$，则有
\begin{equation*}
\frac{D^2}{d\tau^2}S^\mu = R^\mu{}_{\nu\rho\sigma}U^\nu U^\rho S^\sigma.
\tag{7.103}
\end{equation*}
我们希望把右端计算到 $h^{\mathrm{TT}}_{\mu\nu}$ 的一阶。如果取检验粒子缓慢运动，我们可以把四速度表成时间方向的单位矢量加上 $h^{\mathrm{TT}}_{\mu\nu}$ 阶及更高阶的改正；但我们知道 Riemann 张量已经是一阶量，所以 $U^\nu$ 的改正可以忽略，于是写成
\begin{equation*}
U^\nu = (1, 0, 0, 0).
\tag{7.104}
\end{equation*}
因此我们只需计算 $R^\mu{}_{00\sigma}$，或等价地 $R_{\mu 00\sigma}$。由 (7.5) 有
\begin{equation*}
R_{\mu 00\sigma} = \tfrac{1}{2}(\partial_0\partial_0 h^{\mathrm{TT}}_{\mu\sigma} + \partial_\sigma\partial_\mu h^{\mathrm{TT}}_{00} - \partial_\sigma\partial_0 h^{\mathrm{TT}}_{\mu 0} - \partial_\mu\partial_0 h^{\mathrm{TT}}_{\sigma 0}).
\tag{7.105}
\end{equation*}
但 $h^{\mathrm{TT}}_{\mu 0} = 0$，所以
\begin{equation*}
R_{\mu 00\sigma} = \tfrac{1}{2}\partial_0\partial_0 h^{\mathrm{TT}}_{\mu\sigma}.
\tag{7.106}
\end{equation*}
另一方面，对缓慢运动的粒子，在最低阶我们有 $\tau = x^0 = t$，于是测地偏离方程成为
\begin{equation*}
\frac{\partial^2}{\partial t^2}S^\mu = \frac{1}{2}S^\sigma\frac{\partial^2}{\partial t^2}h^{\mathrm{TT}\mu}{}_{\sigma}.
\tag{7.107}
\end{equation*}

对我们的沿 $x^3$ 方向传播的波，这意味着只有 $S^1$ 和 $S^2$ 会受到影响——检验粒子只在垂直于波矢量的方向上被扰动。这一点在电磁学中是大家熟悉的：平面波中的电场和磁场都垂直于波矢量。

我们的波由这两个数刻画，为了后面方便，我们把它们重新命名为：
\begin{gather*}
h_+ = C_{11} \\
h_\times = C_{12},
\tag{7.108}
\end{gather*}

% ---- 书页 297 / p312 ----

于是
\begin{equation*}
C_{\mu\nu} =
\begin{pmatrix}
0 & 0 & 0 & 0 \\
0 & h_+ & h_\times & 0 \\
0 & h_\times & -h_+ & 0 \\
0 & 0 & 0 & 0
\end{pmatrix}.
\tag{7.109}
\end{equation*}
让我们分别考虑二者的效应，先讨论 $h_\times = 0$ 的情形。这时有
\begin{equation*}
\frac{\partial^2}{\partial t^2}S^1 = \frac{1}{2}S^1\frac{\partial^2}{\partial t^2}(h_+e^{ik_\sigma x^\sigma})
\tag{7.110}
\end{equation*}
而
\begin{equation*}
\frac{\partial^2}{\partial t^2}S^2 = -\frac{1}{2}S^2\frac{\partial^2}{\partial t^2}(h_+e^{ik_\sigma x^\sigma}).
\tag{7.111}
\end{equation*}
这两个方程可以立即解出，在最低阶给出
\begin{equation*}
S^1 = \left(1 + \tfrac{1}{2}h_+e^{ik_\sigma x^\sigma}\right)S^{1}{}_{(0)}
\tag{7.112}
\end{equation*}
以及
\begin{equation*}
S^2 = \left(1 - \tfrac{1}{2}h_+e^{ik_\sigma x^\sigma}\right)S^{2}{}_{(0)}.
\tag{7.113}
\end{equation*}

于是，初始时刻沿 $x^1$ 方向分开的粒子将在 $x^1$ 方向上振荡，初始沿 $x^2$ 方向分开的粒子亦然。也就是说，如果我们从 $x$-$y$ 平面上一圈静止的粒子出发，那么当波经过时，它们会以“+”的形状来回弹动，如图 7.4 所示。另一方面，对 $h_+ = 0$ 而 $h_\times \neq 0$ 的情形作同样的分析，将得到解
\begin{equation*}
S^1 = S^{1}{}_{(0)} + \tfrac{1}{2}h_\times e^{ik_\sigma x^\sigma}S^{2}{}_{(0)}
\tag{7.114}
\end{equation*}
以及
\begin{equation*}
S^2 = S^{2}{}_{(0)} + \tfrac{1}{2}h_\times e^{ik_\sigma x^\sigma}S^{1}{}_{(0)}.
\tag{7.115}
\end{equation*}

%FIG{7.4}: {具有 + 偏振的引力波会把检验粒子圆环扭曲成按 + 图样振荡的椭圆。}

% ---- 书页 298 / p313 ----

%FIG{7.5}: {具有 $\times$ 偏振的引力波会把检验粒子圆环扭曲成按 $\times$ 图样振荡的椭圆。}

在这种情形下，粒子圆环会来回弹动，保持“$\times$”的形状，如图 7.5 所示。因此 $h_+$ 和 $h_\times$ 这两个记号的含义应当清楚了。这两个量量度引力波的两种独立的线偏振模式，分别称为“+”偏振与“$\times$”偏振（plus and cross）。如果我们愿意，通过定义
\begin{gather*}
h_R = \frac{1}{\sqrt{2}}(h_+ + ih_\times), \\
h_L = \frac{1}{\sqrt{2}}(h_+ - ih_\times).
\tag{7.116}
\end{gather*}
就可以考虑右旋和左旋的圆偏振模式。

一个纯 $h_R$ 波的效应是使粒子沿右手方向转动，如图 7.6 所示；左旋模式 $h_L$ 与此类似。注意，各个粒子并不绕着环运动；它们只是在小本轮（epicycles）上打转。

我们可以把经典引力波的偏振态与量子化后预期会出现的粒子种类联系起来。量子化场的自旋与该场在空间转动下的变换性质直接相关。电磁场有两个独立的偏振态，由 $x$-$y$ 平面内的矢量描写；等价地说，单一偏振模式在这个平面内转动 $360^{\circ}$ 后保持不变。量子化后，这个理论给出光子——一种无质量的自旋为 1 的粒子。而中微子也是一种无质量粒子，描写它的场在转动 $360^{\circ}$ 下改变符号；它在 $720^{\circ}$ 转动下不变，于是我们说它具有

%FIG{7.6}: {具有 R 偏振的引力波会把检验粒子圆环扭曲成沿右手方向旋转的椭圆。}
